% =============================================================
% Context and Retrieval (case-study chapter)
% Chapter order is set in ../main.tex; the rendered chapter number
% comes from \chapter{} below.
%
% This is one of the planned interleaved case-study chapters.
% Pattern: a section or two of theoretical discussion, then a
% concrete worked example with multiple approaches measured
% side-by-side. The methodology behind these chapters lives in
% manuscript/notes/case_studies.md.

\chapter{Context and Retrieval}
\label{ch:context}

\todo{Chapter stub. Subsection structure laid out below; writing
deferred until Chapters~\ref{ch:mental} through~\ref{ch:toolcalls}
have been fleshed out.}

Chapter~\ref{ch:toolcalls} treated the tool call as the unit of
work. One of the most common things a tool actually does is reach
into a body of text the model was not trained on --- a private
codebase, an internal wiki, a freshly-published documentation
page, a chat history older than the context window --- and bring
back a relevant slice. This is the territory of \emph{retrieval},
and it is where Retrieval-Augmented Generation (RAG) lives. The
chapter has two halves. The first surveys the design space:
embeddings-based RAG, agentic retrieval, graph-based retrieval,
and hybrids of the three. The second is a comparative case study
that runs three of those approaches on the same problem and
reports the numbers, so the choice does not have to be made on
intuition alone.

\section{The Problem}
\label{sec:context-problem}

\todo{One short section. Restate the LLM's blind spot from the
intro: training cutoff, private corpora, queries about state the
model cannot have memorised. Map the design space across three
orthogonal axes.}

\paragraph{Corpus structure.}
\todo{Unstructured text (manuals, blog posts), structured text
(code with grep-able identifiers), structured data (database
rows, knowledge graphs). Each axis end has a different natural
retrieval method.}

\paragraph{Corpus size.}
\todo{Fits-in-context vs does-not. Below the threshold,
``retrieval'' is just paste-the-whole-corpus; above the threshold
the design space opens up. The threshold has moved with model
context windows but the question still applies.}

\paragraph{Query shape.}
\todo{Keyword-precise (the user names a function, an env var, a
ticket number), semantic-fuzzy (the user describes a behaviour
without using the right keywords), multi-hop (``find functions
that call X and also write to the database'').}

\section{Approaches}
\label{sec:context-approaches}

\todo{One subsection per approach. Theoretical discussion only;
the runnable code lives in \S\ref{sec:context-case-study}.}

\subsection{Classic RAG with embeddings}
\label{sec:context-rag}

\todo{Chunk the corpus, embed every chunk, store the embeddings,
embed the query at run time, retrieve the nearest~$k$, paste
into the prompt. Trade-offs: chunking strategy is load-bearing;
the embedding model is the bottleneck on quality; recall on
multi-hop queries is poor; works at scale once vector indexes
are involved. The canonical illustrative reference is
\citet{bendersky2023ragingo}, whose 100-line Go implementation
is the model for the case study below.}

\subsection{Agentic retrieval}
\label{sec:context-agentic}

\todo{Hand the agent a shell or a code interpreter and let it
look. \code{grep}, \code{find}, \code{rg}, \code{curl} for
external sources, \code{git log} for code history. The agent
iterates: a first search rarely returns exactly what is needed,
but the next search is informed by what it just read.
Trade-offs: works on the structured end of the corpus axis
(code, configs, logs); poor on narrative text where the right
keywords do not appear; introduces extra round-trips with the
model; consumes context as it goes
(\S\ref{sec:context}). The shell-tool design of
\S\ref{sec:shell} is the substrate.}

\subsection{Graph-based retrieval}
\label{sec:context-graph}

\todo{Build a graph (call graph, dependency graph, entity-relation
knowledge graph), navigate it. The agent's queries become graph
queries: ``functions that call \code{X}'', ``papers that cite
\code{Y}'', ``people who report into \code{Z}''. Trade-offs: graph
construction is expensive; the schema is design work; multi-hop
queries that frustrate RAG and grep are this approach's
strength. Often the right choice for code repositories, large
structured knowledge bases, and any retrieval task with explicit
relations to traverse. We sketch the design here; the case study
implements it conceptually rather than fully.}

\subsection{Hybrid approaches}
\label{sec:context-hybrid}

\todo{The interesting question is rarely ``which one'' but
``which combination''. A common pattern is:
keyword/grep first (cheap, exact); embeddings on miss (semantic,
fuzzy); graph for the explicit-relation queries (correct,
expensive). The agent is a natural orchestrator for this stack
because it can decide which to try based on what the query looks
like. We work the hybrid up in the case study below.}

\section{Case Study: the Bendersky Problem, Modernised}
\label{sec:context-case-study}

\todo{The substantive section. Three implementations of the same
task, run on the same benchmark, with measurable results.}

\subsection{Setup}
\label{sec:context-setup}

\todo{The original problem from
\citet{bendersky2023ragingo}: answer questions about the Go
documentation that the model was not trained on. Restate as: a
read-only corpus of $\sim$300 Markdown files; a benchmark of
20~hand-written questions covering the three query-shape buckets
(keyword, semantic, multi-hop); a fixed evaluation rubric; a
fixed Gemma~4 / Ollama setup. All inputs and outputs go into
\fname{manuscript/scripts/case-studies/context-and-retrieval/}.}

\subsection{Approach A: Classic RAG}
\label{sec:context-approach-a}

\todo{A faithful Python port of the Bendersky pipeline.
SQLite for chunk storage; OpenAI's
\code{text-embedding-3-small} or a local equivalent for
embeddings; cosine similarity in a loop (no vector DB, since the
corpus is small enough); top-3 chunks pasted into the prompt.
Approximately 100~lines.}

\subsection{Approach B: Agentic Retrieval}
\label{sec:context-approach-b}

\todo{The Chapter~\ref{ch:toolcalls} tool-call loop, with two
tools: \code{shell} (so the agent can grep, find, cat) and
\code{web\_search} (for queries that genuinely need an outside
source). No embeddings, no SQLite. Approximately 50~lines on top
of code we already wrote.}

\subsection{Approach C: Hybrid}
\label{sec:context-approach-c}

\todo{Approach B with the embeddings index of Approach A
exposed as a third tool, \code{semantic\_search}. The agent
chooses freely between \code{shell} and \code{semantic\_search};
the system prompt nudges it toward the cheaper shell tool first.
Approximately 30~more lines.}

\subsection{Approach D (sketch only): Graph-Based}
\label{sec:context-approach-d}

\todo{For the Go documentation corpus, a graph layer is overkill
and will not appear in the case study's measured comparison. We
include a half-page sketch of what the graph would look like
(documents as nodes, with edges for ``links to'' and ``is
referenced by''), what queries it would unlock (``which doc page
do most other doc pages link to about
\code{GOTOOLCHAIN}?''), and a forward reference to the agents
chapters where graph retrieval becomes load-bearing.}

\section{Measuring the Approaches}
\label{sec:context-measure}

\todo{The numbers. Each approach is run on the same 20~questions;
each answer is graded by a fixed rubric (correct / partially
correct / wrong / refusal); the wall-clock time and token cost
are recorded. The expected shape of the table:
Approach~A wins on semantic queries; Approach~B wins on
keyword and exact-match queries and is dramatically cheaper;
Approach~C wins overall by combining them. Numbers will be filled
in once the case-study scripts are written.}

\section{When to Reach for Which}
\label{sec:context-when}

\todo{Recommendations. The structure mirrors
\S\ref{sec:mcp-recs}: a handful of named cases, each with a clear
default and the conditions that flip it.}

\begin{enumerate}[label=\textbf{(\arabic*)},leftmargin=2.2em,itemsep=0.5em]
\item \todo{If the corpus fits in the context window and the cost
  is acceptable, do not retrieve. Paste it.}
\item \todo{If the corpus is structured and the queries are
  keyword-shaped (code, configs, logs), reach for the agentic
  approach first. It is cheaper, simpler, and often more accurate.}
\item \todo{If the corpus is unstructured and the queries are
  semantic, classic RAG remains the right baseline.}
\item \todo{If the queries cross the boundary between keyword and
  semantic --- which most real ones do --- the hybrid is the
  default.}
\item \todo{If the corpus has explicit relations the agent will
  query (call graph, citation graph, entity graph), the graph
  layer earns its keep. Otherwise it is overkill.}
\end{enumerate}

\section{Memory: Retrieval Across Time}
\label{sec:context-memory}

\todo{A short closing section, half a page. Retrieval as treated
above is per-query: the corpus is fixed, the agent looks something
up, the conversation moves on. Long-running agents face a related
but different problem: what to remember from earlier turns of the
\emph{same} conversation, and how to retrieve it later. The same
three approaches apply (embed past turns; let the agent grep its
own scratchpad; build a graph of entities and decisions) with
different trade-offs. We treat memory in detail in the agents
chapters; this section establishes the connection.}
